\section{Programação Dinâmica}



\subsection{Divide and Conquer DP Optimization}

Otimiza DPs da forma $dp[i][j] = \min_{k < j} (dp[i-1][k] + C(k, j))$ quando há monotonia no ponto ótimo ($opt(i, j) \le opt(i, j+1)$).

\textbf{Complexidade:} $O(k n \log n)$.

\begin{lstlisting}
ll dp_before[MAXN], dp_curr[MAXN];
ll cost(int i, int j) {
    // Retorna o custo de transicao do intervalo [i, j]
    return 0; 
}
void compute(int l, int r, int opt_l, int opt_r) {
    if (l > r) return;
    int mid = (l + r) / 2;
    pair<ll, int> best = {1e18, -1};

    for (int k = opt_l; k <= min(mid, opt_r); k++) 
        best = min(best, {dp_before[k] + cost(k + 1, mid), k});

    dp_curr[mid] = best.first;
    int opt = best.second;
    compute(l, mid - 1, opt_l, opt);
    compute(mid + 1, r, opt, opt_r);
}
\end{lstlisting}


\subsection{SOS DP (Sum Over Subsets)}

Calcula a soma sobre todos os subconjuntos de cada máscara de bits.

\textbf{Complexidade:} $O(n 2^n)$.

\begin{lstlisting}
// F[mask] armazena o valor inicial. 
// Ao termino do loop, F[mask] contera a soma dos valores de todos os subconjuntos de mask.
void sos_dp(int n, vector<ll>& F) {
    for (int i = 0; i < n; ++i) {
        for (int mask = 0; mask < (1 << n); ++mask) {
            if (mask & (1 << i)) {
                F[mask] += F[mask ^ (1 << i)];
            }
        }
    }
}
\end{lstlisting}


\subsection{Digit DP Template}

Conta números em um intervalo $[0, n]$ que satisfazem certas propriedades baseadas em dígitos usando memoização.

\begin{lstlisting}
string s;
ll memo[20][2]; // Ajuste os estados extras conforme o problema
ll dp(int idx, bool tight, /* outros estados necessarios */) {
    if (idx == s.size()) return 1; // Retorna 1 se encontrou um numero valido
    if (memo[idx][tight] != -1) return memo[idx][tight];

    ll ans = 0;
    int limit = tight ? (s[idx] - '0') : 9;
    
    for (int d = 0; d <= limit; ++d) {
        // Validacoes extras para o digito 'd' entram aqui se necessario
        ans += dp(idx + 1, tight && (d == limit) /*, novos estados */);
    }
    return memo[idx][tight] = ans;
}

ll solve(ll n) {
    if (n < 0) return 0;
    s = to_string(n);
    memset(memo, -1, sizeof(memo));
    return dp(0, true /*, estado inicial */);
}
// Para resolver [A, B], faca solve(B) - solve(A - 1)
\end{lstlisting}